% BHU PYQ -- Manually transcribed & error-corrected
% Paper: BCH-215 : Fundamentals of Business Finance
% Exam: B.Com. (Hons.) Semester III Examination, 2013-14

\documentclass[11pt,a4paper]{article}
\usepackage[a4paper,top=2.5cm,bottom=2.5cm,left=2.8cm,right=2.8cm,headheight=14pt]{geometry}
\usepackage{amsmath,amssymb}
\usepackage[T1]{fontenc}
\usepackage[utf8]{inputenc}
\usepackage{lmodern,microtype,fancyhdr,enumitem,hyperref,lastpage,parskip}

\hypersetup{
    colorlinks=true,
    urlcolor=black,
    linkcolor=black,
    pdftitle={BCH-215: Fundamentals of Business Finance -- BHU B.Com. (Hons.) Sem III 2013-14}
}

\pagestyle{fancy}
\fancyhf{}
\renewcommand{\headrulewidth}{0.4pt}
\renewcommand{\footrulewidth}{0.4pt}
\fancyhead[L]{\small   \textit{Banaras Hindu University}}
\fancyhead[R]{\small   \textit{BCH-215: Fundamentals of Business Finance}}
\fancyfoot[C]{\small Page  \thepage\ of \pageref{LastPage}}


\newlist{parts}{enumerate}{2}
\setlist[parts,1]{label=(\alph*),leftmargin=*,itemsep=5pt,topsep=3pt}
\setlist[parts,2]{label=(\roman*),leftmargin=*,itemsep=3pt,topsep=2pt}

\newcommand{\pts}[1]{\hfill   \textit{[#1]}}


\begin{document}


\begin{center}
  {\large\bfseries BANARAS HINDU UNIVERSITY}\\[2pt]
  {
\normalsize B.Com. (Hons.) --- Semester III Examination, 2013-14}\\[6pt]
  \rule{0.8\linewidth}{0.5pt}\\[6pt]
  {\large\bfseries Paper No. BCH-215: Fundamentals of Business Finance}\\[4pt]
  {
\normalsize Time: 3 Hours\hfill Maximum Marks: 70}
\end{center}

\rule{\linewidth}{0.4pt}

\smallskip



\noindent   \textit{Instructions: Answer all questions. All questions carry equal marks (14 marks each).}

\bigskip




\noindent   \textbf{Question 1.} \\
Define business finance. State the traditional and modern views of business finance.\pts{14}

\centerline{   \textbf{OR}}

What are the objectives of business finance? Distinguish between the profit maximisation and wealth maximisation.\pts{14}

\medskip\rule{\linewidth}{0.2pt}




\noindent   \textbf{Question 2.} \\
Define financial planning. Discuss the various steps involved in financial planning.\pts{14}

\centerline{   \textbf{OR}}

What are the main features of financial planning? Describe its importance and limitations.\pts{14}

\medskip\rule{\linewidth}{0.2pt}




\noindent   \textbf{Question 3.} \\
What do you mean by overcapitalisation? State the causes, impact and remedies of overcapitalisation.\pts{14}

\centerline{   \textbf{OR}}

Write short notes on the following:\pts{7+7}

\begin{parts}
  \item Undercapitalisation
  \item Fixed capital
\end{parts}

\medskip\rule{\linewidth}{0.2pt}




\noindent   \textbf{Question 4.} \\
State the various sources of long-term capital.\pts{14}

\centerline{   \textbf{OR}}

Write short notes on the following:\pts{7+7}

\begin{parts}
  \item Debenture
  \item Preference share capital
\end{parts}

\medskip\rule{\linewidth}{0.2pt}




\noindent   \textbf{Question 5.} \\
Define working capital. State the various components of working capital and its importance for the business enterprises.\pts{14}

\centerline{   \textbf{OR}}

Write short notes on the following:\pts{7+7}

\begin{parts}
  \item Factoring
  \item Certificate of deposits
\end{parts}

\vfill
\rule{\linewidth}{0.4pt}

\begin{center}\small
\href{https://bhu-exam.vercel.app/}{Click here to visit our website}\\[4pt]
\href{https://me-aryan.vercel.app/}{Click here to visit our contributor}\\[4pt]
\href{https://quashan-me.vercel.app/}{Click here to visit our contributor}
\end{center}

\end{document}
